\documentclass[11pt,reqno]{amsart}
\usepackage{amsmath}
\usepackage{amsthm}
\usepackage{amsfonts}
\usepackage{amssymb}
\usepackage{array}
\usepackage{booktabs}
\usepackage[english]{babel}
\usepackage{placeins}

\frenchspacing
\allowdisplaybreaks
\linespread{1.2}
\usepackage[a4paper, lmargin=0.15\paperwidth, rmargin=0.15\paperwidth,
            tmargin=0.1\paperheight, bmargin=0.1\paperheight]{geometry}
\usepackage[colorlinks = true,
            linkcolor = blue,
            urlcolor  = blue,
            citecolor = blue,
            anchorcolor = blue]{hyperref}

\theoremstyle{plain}
\newtheorem{theorem}{Theorem}[section]
\newtheorem{lemma}[theorem]{Lemma}
\newtheorem{proposition}[theorem]{Proposition}
\newtheorem{corollary}[theorem]{Corollary}

\theoremstyle{definition}
\newtheorem{definition}[theorem]{Definition}

\theoremstyle{remark}
\newtheorem{remark}[theorem]{Remark}

\newcommand{\T}{\mathbb T}
\newcommand{\Z}{\mathbb Z}
\newcommand{\R}{\mathbb R}
\newcommand{\N}{\mathbb N}
\DeclareMathOperator{\dist}{dist}
\DeclareMathOperator{\tr}{tr}
\DeclareMathOperator{\Var}{Var}

\title[Counting Partial Hadamard Matrices in the Near-Quadratic Regime]
{Counting Partial Hadamard Matrices\\in the Near-Quadratic Regime}

\author{Haowei Lin}
\address{Tencent Hunyuan}
\email{linhaowei@pku.edu.cn}

\date{}

\begin{document}

\begin{abstract}
Let \(N_{n,m}\) denote the number of \(n\times m\) matrices with entries in
\(\{\pm1\}\) whose rows are pairwise orthogonal, and suppose that \(m\) is a
multiple of four. We prove an explicit, uniform asymptotic formula in a
near-quadratic range. Writing
\[
  A_{n,m}=2^{nm+(n-1)^2}(2\pi m)^{-n(n-1)/4},\qquad
  R_{n,m}=\exp\Bigl(-\frac{n(n-1)(2n-1)}{24m}\Bigr),
\]
we show that \(N_{n,m}=[1+o(1)]A_{n,m}R_{n,m}\) uniformly whenever
\(m/(n^2\log(2\pi m))\to\infty\). In particular, for every fixed \(\eta>0\) the
formula holds uniformly in the regime \(m/n^{2+\eta}\to\infty\). The correction
factor \(R_{n,m}\) is essential below the cubic scale and resums the leading
cubic-scale correction. Consequently the infimum of the admissible power-law
exponents is at most \(2\), improving the previous bound \(3\); validity at the
endpoint exponent \(2\) is not claimed.

The proof proceeds by adjoining rows one at a time. For a fixed \(k\)-row
partial Hadamard matrix, the number of admissible next rows is a
\(k\)-dimensional Fourier integral. Exact orthogonality identities give a
uniform central expansion for every partial Hadamard matrix. Away from the
universal saddle lattice, a scale-adapted truncation and a Bernstein bound show
exponential decay for all but an exponentially small set of sign matrices. A
quantitative induction then shows that this exceptional set is negligible among
partial Hadamard matrices as soon as \(m\gg n^2\log m\). We also give a
closed-form error bound valid for every integer pair \(n\ge2\) and
\(m\in4\N\).
\end{abstract}

\maketitle

\section{Introduction}

A real partial Hadamard matrix is a rectangular sign matrix whose rows are
pairwise orthogonal. These objects interpolate between elementary orthogonal
sign systems and square Hadamard matrices, and their counting problem is
naturally linked to lattice random walks and Fourier inversion. De Launey and
Levin introduced a Fourier-analytic framework for this problem and obtained
asymptotic counting results for fixed row dimension and, implicitly, in a
high-dimensional range \cite{deLauneyLevin2010}. Subsequent work reduced the
required separation between the width and the number of rows. In particular,
Davis proved a precise expansion at and above the cubic scale \(m\asymp n^3\),
extending an unpublished estimate of Canfield \cite{Canfield2011,Davis2026}.

The standard leading term is
\[
  A_{n,m}=2^{nm+(n-1)^2}(2\pi m)^{-n(n-1)/4}.
\]
At the cubic scale this term alone is no longer asymptotically exact. The first
correction found in \cite{Davis2026} is of order \(n^3/m\). The main point of
the present argument is that the whole correction can be resummed into the
elementary factor
\[
  R_{n,m}=\exp\Bigl(-\frac{n(n-1)(2n-1)}{24m}\Bigr).
\]
Once this factor is retained, the central part of the Fourier integral remains
stable down to the near-quadratic scale. The remaining obstacle is not the
local expansion but the possibility of additional, highly structured far-field
peaks. Such peaks do occur for exceptional partial Hadamard matrices, so a
uniform pointwise bound for every matrix is false; see
Remark~\ref{rem:exceptional}. Instead, we prove that the exceptional matrices
are exponentially rare among all sign matrices, and then compare this rarity
with a quantitative lower bound for the density of partial Hadamard matrices.

Our natural asymptotic condition is
\[
  \frac{m}{n^2\log(2\pi m)}\longrightarrow\infty.
\]
This implies the formula in every power-law regime \(m/n^{2+\eta}\to\infty\)
with fixed \(\eta>0\). It does not settle the endpoint condition
\(m/n^2\to\infty\): for example, \(m=n^2\log\log n\) does not force
\(m/(n^2\log m)\to\infty\). The endpoint remains outside the reach of the
present exceptional-set argument.

\subsection*{Relation to the cubic expansion}
If \(m=4t\) and \(t=\Theta n^3\) with fixed \(\Theta>0\), then
\[
  R_{n,4t}\longrightarrow e^{-1/(48\Theta)}.
\]
For large \(\Theta\) this gives
\(e^{-1/(48\Theta)}=1-\tfrac1{48\Theta}+O(\Theta^{-2})\), which agrees with the
leading correction in \cite{Davis2026}. Thus the present formula is consistent
with the known cubic expansion while also predicting the full limiting
correction factor.

\section{Main results}

Throughout the paper, \(m=4t\ge4\) and
\[
  L:=\log(2\pi m)=\log(8\pi t).
\]
For integers \(n\ge2\) and \(t\ge1\), define
\begin{equation}\label{eq:B}
  B_{n,m}:=2^{nm+(n-1)^2}(2\pi m)^{-n(n-1)/4}
    \exp\Bigl(-\frac{n(n-1)(2n-1)}{24m}\Bigr),\qquad(m=4t),
\end{equation}
\begin{multline}\label{eq:E}
  E(n,t):=\frac{n^4}{4m^2}+2ne^{-m/(4n)}+2ne^{2n-m/(86n)}
    +3ne^{(n/2)L-m/32}\\
    +2n\exp\Bigl(2n^2L+2n\log(200n^2)-\frac m{1024}\Bigr),
\end{multline}
\begin{equation}\label{eq:S}
  S(n,t):=\Bigl(\frac{n^2L}4+\frac{n^3}{12m}\Bigr)
    \Bigl(\frac{20000\,n^2L}{m}\Bigr)^{n^2},
\end{equation}
\begin{equation}\label{eq:eps}
  \varepsilon(n,t):=\exp\bigl(E(n,t)+S(n,t)\bigr)-1.
\end{equation}

\begin{theorem}[Explicit all-parameter estimate]\label{thm:explicit}
For every pair of integers \(n\ge2\) and \(t\ge1\),
\[
  \Bigl|\frac{N_{n,4t}}{B_{n,4t}}-1\Bigr|\le\varepsilon(n,t).
\]
\end{theorem}

The closed form in \eqref{eq:eps} is intentionally conservative. Its purpose is
to make the statement valid without side conditions, including parameter pairs
for which no \(n\times4t\) partial Hadamard matrix can exist.

\begin{theorem}[Near-quadratic uniform asymptotic]\label{thm:uniform}
Let \((n_i,t_i)\) be any sequence with \(n_i\ge2\), \(t_i\ge1\), and put
\(m_i=4t_i\). If
\[
  \frac{m_i}{n_i^2\log(2\pi m_i)}\longrightarrow\infty,
\]
then \(N_{n_i,m_i}/B_{n_i,m_i}\to1\). The convergence is uniform in the
displayed regime.
\end{theorem}

\begin{corollary}[Every exponent strictly larger than two]\label{cor:exponent}
For every fixed \(u>2\),
\[
  N_{n,4t}=[1+o(1)]B_{n,4t}
\]
uniformly whenever \(t/n^u\to\infty\).
\end{corollary}

If a critical exponent is defined as the infimum of the admissible power
exponents, Corollary~\ref{cor:exponent} gives an upper bound of \(2\). This
infimal statement must not be confused with validity at the endpoint \(u=2\),
which is not proved here. Section~\ref{sec:asymptotic} makes the infimum
statement precise, and Section~\ref{sec:limits} records what is not proved.

\section{Fourier representation and row extension}

Write \(\T=\R/(2\pi\Z)\), with total Lebesgue measure \(2\pi\). For a
\(k\times m\) sign matrix \(H\), let \(h_1,\dots,h_k\in\{\pm1\}^m\) be its rows
and \(H_1,\dots,H_m\in\{\pm1\}^k\) its columns. Let \(\mathcal P_k\) be the set
of \(k\)-row partial Hadamard matrices of width \(m\), so that
\(|\mathcal P_k|=N_{k,m}\) and \(N_{1,m}=2^m\).

For \(\theta\in\T^k\), define
\[
  \Phi_H(\theta):=\prod_{C=1}^m\cos(\theta\cdot H_C),
\]
and let \(X(H):=\#\{x\in\{\pm1\}^m:Hx=0\}\). Every \((k+1)\)-row partial
Hadamard matrix is uniquely obtained by adjoining one such row, hence
\begin{equation}\label{eq:recursion}
  N_{k+1,m}=\sum_{H\in\mathcal P_k}X(H).
\end{equation}

\begin{lemma}[Fourier representation]\label{lem:fourier}
For every \(k\times m\) sign matrix \(H\),
\[
  X(H)=2^m(2\pi)^{-k}\int_{\T^k}\Phi_H(\theta)\,d\theta.
\]
\end{lemma}

\begin{proof}
For an integer \(z\), \(\mathbf 1_{\{z=0\}}=(2\pi)^{-1}\int_\T e^{iz\theta}\,
d\theta\). Apply this identity to each coordinate of \(Hx\), sum over
\(x\in\{\pm1\}^m\), and factor the resulting product:
\[
  X(H)=(2\pi)^{-k}\int_{\T^k}\prod_{C=1}^m
    \Bigl(\sum_{x_C=\pm1}e^{ix_C(\theta\cdot H_C)}\Bigr)d\theta
  =2^m(2\pi)^{-k}\int_{\T^k}\Phi_H(\theta)\,d\theta. \qedhere
\]
\end{proof}

\section{The universal saddle lattice}

Define
\[
  \Lambda_k:=\Bigl\{\tfrac\pi2c\ (\mathrm{mod}\ 2\pi):c\in\Z^k,\
    \textstyle\sum_{j=1}^kc_j\equiv0\ (\mathrm{mod}\ 2)\Bigr\}\subset\T^k.
\]
We call \(\Lambda_k\) the universal saddle lattice. It is universal because
every partial Hadamard integrand is exactly periodic under its translations. It
need not contain every point where a particular integrand has modulus one;
exceptional matrices can have additional far-field maxima, as shown in
Remark~\ref{rem:exceptional}.

\begin{lemma}\label{lem:lattice}
One has \(|\Lambda_k|=2^{2k-1}\), and distinct points of \(\Lambda_k\) are at
least \(\pi/\sqrt2\) apart in \(\T^k\).
\end{lemma}

\begin{proof}
Modulo \(2\pi\), take \(c\in\{0,1,2,3\}^k\). Exactly half of the \(4^k\)
classes have even coordinate sum. For the distance statement, let
\(c''\not\equiv0\ (\mathrm{mod}\ 4)\) be the difference of two even-sum
representatives. If some coordinate of \(c''\) is congruent to \(2\) modulo
\(4\), the distance is at least \(\pi\). Otherwise every coordinate is odd or
\(\equiv0\), the number of odd coordinates is positive and even, hence at least
two, giving distance at least \(\sqrt{2(\pi/2)^2}=\pi/\sqrt2\).
\end{proof}

\begin{lemma}[Exact periodicity]\label{lem:periodic}
Let \(H\) be a \(k\)-row partial Hadamard matrix and suppose \(4\mid m\). Then
for every \(\omega\in\Lambda_k\) and \(\beta\in\T^k\),
\(\Phi_H(\omega+\beta)=\Phi_H(\beta)\).
\end{lemma}

\begin{proof}
Write \(\omega=(\pi/2)c\) with \(\sum_jc_j\) even. For each column \(C\),
\(\sum_jc_jH_{jC}\) is even, and therefore \(\omega\cdot H_C\in\pi\Z\). Thus
\[
  \cos((\omega+\beta)\cdot H_C)=\epsilon_C\cos(\beta\cdot H_C),
  \qquad\epsilon_C\in\{\pm1\}.
\]
Let \(s_j=\sum_CH_{jC}\) be the \(j\)th row sum. Then
\(\prod_C\epsilon_C=(-1)^{\frac12\sum_jc_js_j}\). It remains to prove
\(\sum_jc_js_j\equiv0\ (\mathrm{mod}\ 4)\). Every \(s_j\) is even. Moreover, if
\(j\ne\ell\) and the two rows are orthogonal, then
\begin{equation}\label{eq:rowsum}
  s_j+s_\ell\equiv0\ (\mathrm{mod}\ 4).
\end{equation}
Indeed, if \(a,b,c,d\) count the four sign patterns of the two rows, then
orthogonality gives \(a+d=m/2\), which is even, and \(s_j+s_\ell=2(a-d)\) is
divisible by four. Let \(T=\{j:c_j\text{ is odd}\}\). Since \(\sum_jc_j\) is
even, \(|T|\) is even. Terms with even \(c_j\) vanish modulo four, while for
\(j\in T\) one has \(c_js_j\equiv s_j\ (\mathrm{mod}\ 4)\). Pair the indices in
\(T\) and apply \eqref{eq:rowsum}. Hence the product of the signs is \(1\).
\end{proof}

Put \(\rho_k=k^{-1/2}\). By Lemma~\ref{lem:lattice}, the closed balls
\(B(\omega,\rho_k)\) are pairwise disjoint. Define
\[
  \mathrm{Near}_k:=\bigcup_{\omega\in\Lambda_k}B(\omega,\rho_k),\qquad
  \mathrm{Far}_k:=\T^k\setminus\mathrm{Near}_k.
\]
For a \(k\)-row partial Hadamard matrix set
\[
  I(H):=\int_{\|\beta\|\le\rho_k}\Phi_H(\beta)\,d\beta,\qquad
  J(H):=\int_{\mathrm{Far}_k}\Phi_H(\theta)\,d\theta.
\]
Then Lemmas~\ref{lem:fourier} and~\ref{lem:periodic} give
\begin{equation}\label{eq:split}
  X(H)=2^m(2\pi)^{-k}\bigl[2^{2k-1}I(H)+J(H)\bigr].
\end{equation}

\section{The central region}

Throughout this section \(H\) is a \(k\)-row partial Hadamard matrix and
\(y_C=\beta\cdot H_C\).

\begin{lemma}[Exact orthogonality identities]\label{lem:orth}
For every \(\beta\in\R^k\),
\[
  \sum_{C=1}^m(\beta\cdot H_C)^2=m\|\beta\|^2,\qquad
  \|H_C\|^2=k,\qquad
  \sum_{C,C'=1}^m(H_C\cdot H_{C'})^2=m^2k.
\]
\end{lemma}

\begin{proof}
The partial Hadamard condition is exactly \(HH^{\mathsf T}=mI_k\). The first
identity is \(\beta^{\mathsf T}HH^{\mathsf T}\beta=m\|\beta\|^2\), and the
second is immediate. For the third,
\(\sum_{C,C'}(H_C\cdot H_{C'})^2=\tr((HH^{\mathsf T})^2)=m^2k\).
\end{proof}

\begin{lemma}[Taylor bound]\label{lem:taylor}
For \(|y|\le1\),
\[
  -\log\cos y=\frac{y^2}2+\frac{y^4}{12}+s(y),\qquad
  0\le s(y)\le\frac{y^6}{30}.
\]
\end{lemma}

\begin{proof}
The power series of \(-\log\cos y\) on \(|y|<\pi/2\) has positive even
coefficients, beginning with \(y^2/2+y^4/12\). Hence \(s(y)/y^6\) is
nondecreasing in \(|y|\) on \([0,1]\). The numerical inequality
\(-\log(\cos1)-\tfrac7{12}<\tfrac1{30}\) then gives the claim. This inequality
may be verified, for example, by alternating-series bounds for \(\cos1\) and a
positive-term truncation of the exponential series.
\end{proof}

Let \(\gamma\) denote the Gaussian law \(N(0,I_k/m)\) on \(\R^k\), and define
\(Q(\beta):=\sum_{C=1}^my_C^4\) and \(Y_6(\beta):=\sum_{C=1}^my_C^6\).

\begin{lemma}[Gaussian moments]\label{lem:moments}
One has
\[
  \mathbb E_\gamma Q=\frac{3k^2}m,\qquad
  \Var_\gamma(Q)\le\frac{96k^3}{m^2},\qquad
  \mathbb E_\gamma Y_6=\frac{15k^3}{m^2}.
\]
\end{lemma}

\begin{proof}
The vector \((y_1,\dots,y_m)\) is centered Gaussian with
\(\Var(y_C)=k/m\) and \(\mathrm{Cov}(y_C,y_{C'})=(H_C\cdot H_{C'})/m\). The
first and third identities follow from the fourth and sixth moments of a
centered Gaussian. For the variance, if \((Y,Z)\) is centered jointly Gaussian
with variances \(a,b\) and covariance \(r\), Wick's formula gives
\(\mathbb E[Y^4Z^4]=9a^2b^2+72abr^2+24r^4\). Summing this identity and using
Lemma~\ref{lem:orth} yields
\[
  \mathbb EQ^2\le\frac{9k^4}{m^2}+\frac{72k^3}{m^2}+\frac{24k^3}{m^2}.
\]
Subtracting \((\mathbb EQ)^2=9k^4/m^2\) proves the variance bound.
\end{proof}

\begin{lemma}[Chi-square tail]\label{lem:chi2}
If \(Z\sim\chi^2_k\) and \(\lambda\ge1\), then
\(\mathbb P(Z\ge\lambda k)\le\exp\bigl(-\tfrac k2(\lambda-1-\log\lambda)\bigr)\).
If \(\lambda\ge6\), then \(\mathbb P(Z\ge\lambda k)\le e^{-k\lambda/4}\).
\end{lemma}

\begin{proof}
Apply the exponential Markov inequality and optimize the moment generating
function \((1-2s)^{-k/2}\) at \(s=(1-1/\lambda)/2\). The second claim follows
from \(\lambda-1-\log\lambda\ge\lambda/2\) for \(\lambda\ge6\).
\end{proof}

\begin{proposition}[Uniform central contribution]\label{prop:central}
Assume \(m\ge6k^2\). For every \(k\)-row partial Hadamard matrix \(H\),
\[
  I(H)=\Bigl(\frac{2\pi}m\Bigr)^{k/2}e^{-k^2/(4m)}(1+\eta_H),\qquad
  |\eta_H|\le\frac{k^3}{m^2}+2e^{-m/(4k)}.
\]
\end{proposition}

\begin{proof}
On \(A=\{\|\beta\|\le k^{-1/2}\}\), Cauchy--Schwarz gives \(|y_C|\le1\). By
Lemmas~\ref{lem:orth} and~\ref{lem:taylor},
\[
  \Phi_H(\beta)=\exp\Bigl(-\frac{m\|\beta\|^2}2-\frac{Q(\beta)}{12}
    -S_6(\beta)\Bigr),\qquad 0\le S_6(\beta)\le\frac{Y_6(\beta)}{30}.
\]
Therefore \(I(H)=(2\pi/m)^{k/2}\mathbb E_\gamma[\mathbf 1_Ae^{-Q/12-S_6}]\).
Put \(\mu=\mathbb EQ/12=k^2/(4m)\le1/24\) and \(W=Q-\mathbb EQ\). Since
\(W/12\ge-\mu\), Taylor's theorem gives
\[
  \mathbb Ee^{-W/12}\le1+e^\mu\frac{\Var(Q)}{288}\le1+0.35\frac{k^3}{m^2}.
\]
This yields the required upper bound. For the lower bound, Jensen's inequality
gives \(\mathbb Ee^{-Q/12}\ge e^{-\mu}\). Also
\(\mathbb E[\mathbf 1_AS_6]\le\tfrac1{30}\mathbb EY_6=k^3/(2m^2)\). Finally
\(m\|\beta\|^2\sim\chi^2_k\), and \(m/k^2\ge6\), so Lemma~\ref{lem:chi2} gives
\(\mathbb P(A^c)\le e^{-m/(4k)}\). Combining these estimates and using
\(e^{1/24}<1.043\) gives
\[
  \mathbb E[\mathbf 1_Ae^{-Q/12-S_6}]\ge e^{-\mu}
    \Bigl(1-0.53\frac{k^3}{m^2}-1.05e^{-m/(4k)}\Bigr).
\]
The stated symmetric bound follows after harmless rounding.
\end{proof}

\section{The far field}

The estimates in this section apply first to arbitrary sign matrices, without
orthogonality assumptions. For \(0<\tau\le2\), define
\[
  g_\tau(y):=\min\{1-\cos(2y),\tau\},\qquad
  G_\tau(\theta,H):=\sum_{C=1}^mg_\tau(\theta\cdot H_C).
\]

\begin{lemma}[Product bound]\label{lem:product}
For every \(\theta\in\T^k\),
\(|\Phi_H(\theta)|\le\exp\bigl(-\tfrac14G_\tau(\theta,H)\bigr)\).
\end{lemma}

\begin{proof}
Since \(\cos^2z=1-\tfrac{1-\cos2z}2\le1-\tfrac{g_\tau(z)}2\), we have
\(\Phi_H(\theta)^2\le\prod_C\bigl(1-\tfrac{g_\tau(\theta\cdot H_C)}2\bigr)
\le e^{-G_\tau(\theta,H)/2}\). Take square roots.
\end{proof}

For \(\theta\in\T^k\), let
\[
  \mu(\theta):=1-\prod_{j=1}^k\cos(2\theta_j),\qquad
  \sigma(\theta)^2:=\sum_{j=1}^k\dist(\theta_j,(\pi/2)\Z)^2,\qquad
  r(\theta):=\dist(\theta,\Lambda_k).
\]

\begin{lemma}[Geometry of the mean]\label{lem:geometry}
Let \(v\) be uniform on \(\{\pm1\}^k\). Then
\begin{enumerate}
\item \(\mathbb E_v[1-\cos(2\theta\cdot v)]=\mu(\theta)\);
\item \(\mu(\theta)\ge\min\{1/2,r(\theta)^2\}\);
\item if \(\mu(\theta)<1/2\), the nearest point of the full grid
  \((\pi/2)\Z^k\) is unique, belongs to \(\Lambda_k\), and has distance
  \(\sigma(\theta)<1/\sqrt2\).
\end{enumerate}
\end{lemma}

\begin{proof}
The first identity follows from
\(\mathbb E_ve^{2i\theta\cdot v}=\prod_j\cos(2\theta_j)\). Choose a nearest
grid point \(\omega=(\pi/2)c\) and write \(\theta=\omega+\beta\) with
\(|\beta_j|\le\pi/4\), so that \(\sigma(\theta)=\|\beta\|\). If
\(\sum_jc_j\) is odd, then \(\prod_j\cos(2\theta_j)\le0\) and
\(\mu(\theta)\ge1\ge\min\{1/2,r(\theta)^2\}\).

Suppose instead that \(\sum_jc_j\) is even, so that \(\omega\in\Lambda_k\). We
claim \(\sigma(\theta)=r(\theta)\). Indeed \(\Lambda_k\subset(\pi/2)\Z^k\), and
the distance to a subset is at least the distance to the ambient set, so
\(\sigma(\theta)\le r(\theta)\); conversely \(\omega\in\Lambda_k\) gives
\(r(\theta)\le\|\theta-\omega\|=\sigma(\theta)\). (The inclusion alone would
give the inequality in the wrong direction, which is why the parity case
matters here.) Now \(\cos(2\beta_j)\le e^{-2\beta_j^2}\) for
\(|\beta_j|\le\pi/4\), so \(\prod_j\cos(2\beta_j)\le e^{-2\sigma(\theta)^2}\)
and
\[
  \mu(\theta)\ge1-e^{-2\sigma(\theta)^2}
  \ge\min\{1/2,\sigma(\theta)^2\}=\min\{1/2,r(\theta)^2\},
\]
where the middle step holds because \(1-e^{-2s}\ge s\) on \([0,1/2]\) and
\(1-e^{-2s}\ge1-e^{-1}>1/2\) for \(s\ge1/2\).

If \(\mu<1/2\), no coordinate can lie halfway between two grid points, since
that would give \(\cos(2\theta_j)=0\) and hence \(\mu=1\); so the nearest grid
point is unique. The odd-parity case is impossible, and the preceding bound
forces \(\sigma^2<1/2\).
\end{proof}

\begin{lemma}[A truncated mean on dyadic shells]\label{lem:shells}
Suppose \(\theta=\omega+\beta\), where \(\omega\in\Lambda_k\),
\(|\beta_j|\le\pi/4\), and \(\sigma=\|\beta\|\). Let \(p\ge1\) satisfy
\(2^{-p-1}<\sigma^2\le2^{-p}\le\tfrac12\), and put
\(\tau_p=\min\{2,48\cdot2^{-p}\}\). Then
\[
  W_p(\theta):=\mathbb E_vg_{\tau_p}(\theta\cdot v)\ge\tfrac382^{-p},
  \qquad \frac{W_p(\theta)}{\tau_p}\ge\frac1{128}.
\]
\end{lemma}

\begin{proof}
Set \(A(v)=1-\cos(2\theta\cdot v)\). Since \(\sum_jc_j\) is even and
\(v\in\{\pm1\}^k\), we have \(c\cdot v\equiv\sum_jc_j\equiv0\pmod2\), so
\(\omega\cdot v\in\pi\Z\) and \(A(v)=1-\cos(2\beta\cdot v)\): translation by
\(\omega\in\Lambda_k\) does not change \(A\). By Lemma~\ref{lem:geometry} and
\(\sigma^2\le1/2\), \(\mathbb EA\ge\sigma^2\). Moreover
\(1-\cos2x=2\sin^2x\le2x^2\) gives \(A\le2(\beta\cdot v)^2\), and the fourth
moment of a Rademacher sum, \(\mathbb E(\beta\cdot v)^4\le3\sigma^4\), gives
\(\mathbb EA^2\le4\mathbb E(\beta\cdot v)^4\le12\sigma^4\).
If \(p\le4\), then \(48\cdot2^{-p}\ge3>2\), so \(\tau_p=2\), the truncation is
inactive, and the conclusion follows from
\(W_p=\mathbb EA>2^{-p-1}\ge\tfrac382^{-p}\). If \(p\ge5\), then using
\(\sigma^2\le2^{-p}\) and \(\mathbb EA\ge\sigma^2\),
\[
  W_p\ge\mathbb EA-\mathbb E[A\mathbf 1_{\{A>\tau_p\}}]
  \ge\mathbb EA-\frac{\mathbb EA^2}{\tau_p}
  \ge\mathbb EA-\frac{12\sigma^4}{48\cdot2^{-p}}
  \ge\mathbb EA-\frac{\sigma^2}4
  \ge\frac34\mathbb EA>\frac382^{-p}.
\]
Dividing by \(\tau_p\) gives the second bound, since
\(\tfrac38/48=\tfrac1{128}\).
\end{proof}

\begin{lemma}[Lipschitz estimate]\label{lem:lipschitz}
For all \(\theta,\theta'\in\T^k\),
\(|G_\tau(\theta,H)-G_\tau(\theta',H)|\le2m\sqrt k\,\|\theta-\theta'\|\).
\end{lemma}

\begin{proof}
The function \(g_\tau\) is \(2\)-Lipschitz, while
\(|\langle\theta-\theta',H_C\rangle|\le\sqrt k\|\theta-\theta'\|\). Sum over
\(C\).
\end{proof}

\subsection{Shells and nets}
Put \(P:=\max\{0,\lceil\log_2k\rceil-1\}\). Decompose \(\mathrm{Far}_k\) into
\begin{align*}
  F_A&:=\{\theta\in\mathrm{Far}_k:\mu(\theta)\ge1/2\},\\
  F_p&:=\{\theta\in\mathrm{Far}_k:\mu(\theta)<1/2,\
    2^{-p-1}<\sigma(\theta)^2\le2^{-p}\},\qquad1\le p\le P.
\end{align*}
By Lemma~\ref{lem:geometry}, these sets cover \(\mathrm{Far}_k\). Set
\(\eta_A=\tfrac1{16\sqrt k}\) and \(\eta_p=\tfrac{3\cdot2^{-p}}{64\sqrt k}\),
and choose nets \(\mathcal N_A\subset F_A\) and \(\mathcal N_p\subset F_p\)
that are maximal subject to having all pairwise distances \emph{greater than}
the corresponding scale. Maximality then gives covering radius at most that
scale, which is what Proposition~\ref{prop:pointwise} uses.

\begin{lemma}[Net cardinalities]\label{lem:nets}
For all relevant \(p\),
\(|\mathcal N_A|\le e^{k\log(200k^2)}\) and
\(|\mathcal N_p|\le e^{k\log(200k^2)}\).
\end{lemma}

\begin{proof}
Partition \(\T^k\) into axis-parallel cells of side \(h\). A cell has diameter
\(h\sqrt k\), so a set with pairwise distances \emph{greater than} \(\eta\) has
at most one point in each cell when \(h=\eta/\sqrt k\). This gives
\(|\mathcal N_A|\le(1+2\pi\sqrt k/\eta_A)^k=(1+32\pi k)^k\le(102k)^k\),
using \(32\pi<100.54\).

For \(F_p\), every point lies in one of the \(2^{2k-1}\) cubes
\(\omega+[-\pi/4,\pi/4]^k\), \(\omega\in\Lambda_k\), by
Lemma~\ref{lem:geometry}(c). An interval of length \(\pi/2\) on the circle
intersects at most \(\pi/(2h)+2\) cells of side \(h\). Here
\(h=\eta_p/\sqrt k\) and, whenever \(F_p\) is present, \(k\ge3\) and
\(2^p\le k\). Hence
\[
  \frac\pi{2h}+2=\frac{32\pi}3k2^p+2\le34k^2,
\]
since \(32\pi/3<33.52\) and \(0.48k^2\ge2\) for \(k\ge3\). Consequently
\(|\mathcal N_p|\le2^{2k-1}(34k^2)^k\le(136k^2)^k\le(200k^2)^k\). This
corrected endpoint count is the reason for the harmless numerical constant
\(200\).
\end{proof}

\begin{lemma}[Bernstein-type lower tail]\label{lem:bernstein}
Let \(X_1,\dots,X_m\) be i.i.d.\ random variables in \([0,\tau]\) with mean
\(W>0\). Then
\[
  \mathbb P\Bigl(\sum_{C=1}^mX_C\le\frac{mW}2\Bigr)
  \le\exp\Bigl(-\frac{mW}{8\tau}\Bigr).
\]
\end{lemma}

\begin{proof}
Take \(\lambda=1/(2\tau)\). Since \(e^{-z}\le1-z+z^2/2\) for \(z\ge0\) and
\(X_C^2\le\tau X_C\),
\(\mathbb Ee^{-\lambda X_C}\le1-\lambda W+\tfrac{\lambda^2\tau W}2
\le e^{-3\lambda W/4}\). Markov's inequality applied to
\(e^{-\lambda\sum X_C}\) gives the result.
\end{proof}

A uniformly random \(k\times m\) sign matrix has independent uniform columns in
\(\{\pm1\}^k\).

\begin{definition}[Good matrices]\label{def:good}
A \(k\times m\) sign matrix \(H\) is \emph{good} if
\[
  G_2(\theta',H)\ge m/4\quad\text{for every }\theta'\in\mathcal N_A,\qquad
  G_{\tau_p}(\theta',H)\ge(3/16)m2^{-p}\quad\text{for every }
    \theta'\in\mathcal N_p,\ 1\le p\le P.
\]
\end{definition}

\begin{lemma}[Good matrices are typical]\label{lem:typical}
For a uniform random \(k\times m\) sign matrix,
\[
  \mathbb P(H\text{ is not good})\le
    \exp\Bigl(k\log(200k^2)+\log(P+1)-\frac m{1024}\Bigr).
\]
\end{lemma}

\begin{proof}
At a point of \(\mathcal N_A\subset F_A\), the truncation at \(\tau=2\) is
inactive, so \(W=\mu(\theta')\ge1/2\) by the definition of \(F_A\); apply
Lemma~\ref{lem:bernstein} to get failure probability at most \(e^{-m/32}\). At
a point of \(\mathcal N_p\), Lemmas~\ref{lem:shells} and~\ref{lem:bernstein}
give failure probability at most \(e^{-m/1024}\). A union bound over the at
most \(P+1\) nets, each of size at most \(e^{k\log(200k^2)}\) by
Lemma~\ref{lem:nets}, with the weaker per-point bound \(e^{-m/1024}\), proves
the claim. For \(k\le2\) one has \(P=0\) and \(\log(P+1)=0\); the bound remains
valid, and in that range \(\mathrm{Far}_k=F_A\), since a point with
\(\mu<1/2\) would need \(\sigma^2=r^2>1/k\ge1/2\) while \(\mu<1/2\) forces
\(\sigma^2<1/2\).
\end{proof}

\begin{proposition}[Pointwise far-field decay]\label{prop:pointwise}
If \(H\) is good, then
\[
  |\Phi_H(\theta)|\le\exp\Bigl(-\frac{3m}{128}\sigma(\theta)^2\Bigr),
    \quad\theta\in F_p,\qquad
  |\Phi_H(\theta)|\le e^{-m/32},\quad\theta\in F_A.
\]
\end{proposition}

\begin{proof}
For \(\theta\in F_p\), choose \(\theta'\in\mathcal N_p\) with
\(\|\theta-\theta'\|\le\eta_p\). By Definition~\ref{def:good} and
Lemma~\ref{lem:lipschitz},
\[
  G_{\tau_p}(\theta,H)\ge\frac{3m}{16}2^{-p}-2m\sqrt k\eta_p
  =\frac{3m}{32}2^{-p}.
\]
Apply Lemma~\ref{lem:product} and use \(\sigma(\theta)^2\le2^{-p}\). The
argument on \(F_A\) is identical:
\(G_2(\theta,H)\ge m/4-2m\sqrt k\eta_A=m/8\).
\end{proof}

\begin{corollary}[Far integral]\label{cor:far}
If \(H\) is good and \(m\ge128k^2\), then
\[
  |J(H)|\le2^{2k-1}\Bigl(\frac{2\pi}m\Bigr)^{k/2}
    \exp\Bigl(1.54k-\frac{3m}{256k}\Bigr)+(2\pi)^ke^{-m/32}.
\]
\end{corollary}

\begin{proof}
The contribution of \(F_A\) is at most \((2\pi)^ke^{-m/32}\). On the union of
the dyadic shells, use the unique nearest \(\omega\in\Lambda_k\) and
Proposition~\ref{prop:pointwise} to obtain
\[
  \sum_p\int_{F_p}|\Phi_H|\le|\Lambda_k|
    \int_{\|\beta\|>k^{-1/2}}e^{-3m\|\beta\|^2/128}\,d\beta.
\]
The integral is the normalizing constant for a Gaussian with covariance
\(64I_k/(3m)\), multiplied by the tail probability
\(\mathbb P(\chi^2_k>3m/(64k))\). Writing the threshold as \(\lambda k\) gives
\(\lambda=3m/(64k^2)\ge6\). Apply Lemma~\ref{lem:chi2}. Finally
\((64/3)^{k/2}\le e^{1.54k}\).
\end{proof}

\begin{remark}[Why an exceptional set is necessary]\label{rem:exceptional}
There is no uniform estimate \(|\Phi_H(\theta)|\le e^{-cm}\) on all of
\(\mathrm{Far}_k\) for every partial Hadamard matrix. Take \(k=4\),
\(8\mid m\), and use each of the eight vectors \(v\in\{\pm1\}^4\) with
\(v_1v_2v_3v_4=1\) exactly \(m/8\) times as a column. The resulting matrix is a
\(4\)-row partial Hadamard matrix. At \(\theta^*=(\pi/4)(1,1,1,1)\), every
column satisfies \(\theta^*\cdot v\in\pi\Z\), so \(|\Phi_H(\theta^*)|=1\). Yet
\(\dist(\theta^*,\Lambda_4)=\pi/2>1/2=\rho_4\), so
\(\theta^*\in\mathrm{Far}_4\).
\end{remark}

\section{One inductive step}

Fix \(n\ge2\) and \(1\le k\le n-1\). Assume
\begin{equation}\label{eq:regime}
  m\ge20000\,n^2L.
\end{equation}
Then \(m\ge64000n^2\), so all hypotheses of the central and far-field estimates
hold. Define
\begin{equation}\label{eq:V}
  V_k:=2^m2^{2k-1}(2\pi m)^{-k/2}e^{-k^2/(4m)}.
\end{equation}
Since \(2\pi m>16\), one has \(0<V_k\le2^m\).

\begin{proposition}[Extension count for a good partial Hadamard matrix]
\label{prop:extension}
If \(H\) is a good \(k\)-row partial Hadamard matrix, then
\(\bigl|\tfrac{X(H)}{V_k}-1\bigr|\le\delta_k\), where
\[
  \delta_k:=\frac{k^3}{m^2}+2e^{-m/(4k)}
    +1.01e^{1.54k-3m/(256k)}+2.02e^{(k/2)L-m/32}.
\]
\end{proposition}

\begin{proof}
Insert Proposition~\ref{prop:central} and Corollary~\ref{cor:far} into
\eqref{eq:split}. The central main term is exactly \(V_k\). After division by
\(V_k\), the first far-field term is multiplied by \(e^{k^2/(4m)}\le1.01\). The
second is bounded by
\(2e^{kL/2}e^{k^2/(4m)}e^{-m/32}\le2.02e^{kL/2-m/32}\).
\end{proof}

Set \(M_j:=2^m\prod_{i=1}^{j-1}V_i\), with \(M_1=2^m\). Also define
\[
  \psi_k:=2\exp\Bigl(2n^2L+2n\log(200n^2)-\frac m{1024}\Bigr),\qquad
  e_k:=\delta_k+\psi_k.
\]

\begin{proposition}[One-step induction]\label{prop:onestep}
Assume \eqref{eq:regime} and \(N_{k,m}\ge M_k/2\). Then
\[
  \Bigl|\frac{N_{k+1,m}}{N_{k,m}V_k}-1\Bigr|\le e_k.
\]
\end{proposition}

\begin{proof}
Let \(B_k\) be the number of bad \(k\)-row partial Hadamard matrices. By
Lemma~\ref{lem:typical}, the number of bad \(k\times m\) sign matrices is at
most \(2^{km}q_k\), where
\(q_k:=\exp\bigl(2k\log(200k^2)-\tfrac m{1024}\bigr)\), so
\(B_k\le2^{km}q_k\). From \eqref{eq:recursion},
Proposition~\ref{prop:extension}, and \(X(H),V_k\le2^m\),
\[
  \Bigl|\frac{N_{k+1,m}}{N_{k,m}V_k}-1\Bigr|
  \le\delta_k+\frac{B_k2^m}{N_{k,m}V_k}.
\]
It remains to bound the exceptional contribution. Put
\(\Pi_k:=M_k/2^{km}=\prod_{i=1}^{k-1}V_i/2^m\). Since
\(V_i/2^m=2^{2i-1}(2\pi m)^{-i/2}e^{-i^2/(4m)}\), we obtain
\[
  \Pi_k\ge\exp\Bigl(-\frac{k^2L}4-\frac{k^3}{12m}\Bigr)\ge e^{-0.28k^2L}.
\]
Similarly \(V_k/2^m\ge e^{-kL}\). Therefore
\[
  \frac{N_{k,m}V_k}{2^{(k+1)m}}\ge\frac12e^{-0.28k^2L-kL}
  \ge\frac12e^{-2k^2L}.
\]
Consequently
\[
  \frac{B_k2^m}{N_{k,m}V_k}
  \le2\exp\Bigl(2k^2L+2k\log(200k^2)-\frac m{1024}\Bigr)\le\psi_k. \qedhere
\]
\end{proof}

\begin{lemma}[Summed one-step error]\label{lem:summed}
Under \eqref{eq:regime}, \(\sum_{k=1}^{n-1}e_k\le E(n,t)\).
\end{lemma}

\begin{proof}
Each family of summands is increasing in \(k\). Thus
\begin{gather*}
  \sum_{k<n}\frac{k^3}{m^2}\le\frac{n^4}{4m^2},\qquad
  \sum_{k<n}2e^{-m/(4k)}\le2ne^{-m/(4n)},\\
  \sum_{k<n}1.01e^{1.54k-3m/(256k)}\le2ne^{2n-m/(86n)},\qquad
  \sum_{k<n}2.02e^{(k/2)L-m/32}\le3ne^{(n/2)L-m/32},\\
  \sum_{k<n}\psi_k\le2n\exp\Bigl(2n^2L+2n\log(200n^2)-\frac m{1024}\Bigr).
\end{gather*}
Adding the bounds gives \eqref{eq:E}.
\end{proof}

\section{Elementary parameter bounds}

Put \(\Theta:=\tfrac m{n^2L}\). Condition \eqref{eq:regime} is
\(\Theta\ge20000\).

\begin{lemma}[Control of \(E\)]\label{lem:controlE}
If \(\Theta\ge20000\), then \(E(n,t)\le F(\Theta)\), where
\[
  F(x):=\frac1x+2e^{-0.4x}+82e^{-0.009x}+23e^{-0.05x}+471e^{-0.0015x}.
\]
The function \(F\) is decreasing, tends to zero, and satisfies
\(F(20000)<10^{-4}<\tfrac12\).
\end{lemma}

\begin{proof}
Use \(m=n^2L\Theta\), \(L>3.2\), and the elementary estimate
\begin{equation}\label{eq:elem}
  xe^{-axy}\le\frac2{ea}e^{-ay/2},\qquad x,y\ge1,
\end{equation}
which follows from \(xy\ge(x+y)/2\) and optimization in \(x\). The five terms
of \(E\) are bounded respectively by
\(\tfrac1\Theta\), \(2e^{-0.4\Theta}\), \(82e^{-0.009\Theta}\),
\(23e^{-0.05\Theta}\), \(471e^{-0.0015\Theta}\). For the final term, note that
under \(\Theta\ge20000\), \(\log(200n^2)\le3.66L\) and hence
\(2n^2L+2n\log(200n^2)-m/1024\le-m/2048\). The remaining estimates are direct
substitutions into \eqref{eq:elem}. Monotonicity is immediate, and the
numerical value at \(20000\) is less than \(5.001\times10^{-5}\).
\end{proof}

\begin{lemma}[Control of the safety term]\label{lem:controlS}
If \(\Theta\ge40000\), then \(S(n,t)\le G(\Theta)\), where
\(G(x):=\bigl(\tfrac{20000}x\bigr)^4(13+4\log x)\). The function \(G\) is
decreasing on \([40000,\infty)\) and tends to zero.
\end{lemma}

\begin{proof}
Let \(x=20000/\Theta\le1/2\). The prefactor in \eqref{eq:S} is at most
\(n^2L/2\). Since
\(L=\log(2\pi)+2\log n+\log L+\log\Theta\le3.68+4\log n+2\log\Theta\), we get
\[
  S(n,t)\le\frac{n^2}2(3.68+4\log n)x^{n^2}+n^2(\log\Theta)x^{n^2}.
\]
Because \(x\le1/2\) and \(n^2\ge4\), \(x^{n^2}\le2^{4-n^2}x^4\). Elementary
maximization over integers \(n\ge2\) gives
\(n^2(3.68+4\log n)2^{4-n^2}<25.89\) and \(n^22^{4-n^2}\le4\). Thus
\(S\le x^4(13+4\log\Theta)\). Differentiation proves monotonicity.
\end{proof}

\section{Proof of the explicit theorem}

\begin{proof}[Proof of Theorem~\ref{thm:explicit}]
First compute the product of the row factors. Using
\[
  \sum_{k=1}^{n-1}(2k-1)=(n-1)^2,\qquad
  \sum_{k=1}^{n-1}k=\frac{n(n-1)}2,\qquad
  \sum_{k=1}^{n-1}k^2=\frac{n(n-1)(2n-1)}6,
\]
we find
\[
  M_n=2^m\prod_{k=1}^{n-1}V_k
  =2^{mn+(n-1)^2}(2\pi m)^{-n(n-1)/4}
    \exp\Bigl(-\frac{n(n-1)(2n-1)}{24m}\Bigr)=B_{n,4t}.
\]
Suppose first that \eqref{eq:regime} fails. Then \(20000n^2L/m>1\), and hence
\(S(n,t)>\tfrac{n^2L}4+\tfrac{n^3}{12m}\). The trivial bound
\(N_{n,m}\le2^{nm}\) gives
\[
  \frac{N_{n,m}}{B_{n,m}}\le2^{-(n-1)^2}(2\pi m)^{n(n-1)/4}
    \exp\Bigl(\frac{n(n-1)(2n-1)}{24m}\Bigr)
  \le\exp\Bigl(\frac{n^2L}4+\frac{n^3}{12m}\Bigr)<e^{S(n,t)}
  \le1+\varepsilon(n,t).
\]
On the other side, \(N_{n,m}/B_{n,m}\ge0\), so
\(1-N_{n,m}/B_{n,m}\le1\). Since \(S(n,t)>3.2\),
\(\varepsilon(n,t)>e^{3.2}-1>1\). The theorem follows in this case.

Now assume \eqref{eq:regime}. By Lemma~\ref{lem:controlE},
\(E(n,t)<1/2\). We prove inductively that
\[
  N_{k,m}=M_k\prod_{j=1}^{k-1}(1+\xi_j),\qquad|\xi_j|\le e_j,\qquad
  \prod_{j=1}^{k-1}(1+\xi_j)\ge\frac12.
\]
The statement is trivial at \(k=1\). If it holds at \(k\), then
\(N_{k,m}\ge M_k/2\), so Proposition~\ref{prop:onestep} defines \(\xi_k\) with
\(|\xi_k|\le e_k\). By Lemma~\ref{lem:summed},
\(\sum_{j=1}^{n-1}e_j\le E(n,t)\le\tfrac12\). Therefore
\[
  \prod_{j=1}^k(1+\xi_j)\ge\prod_{j=1}^k(1-e_j)
  \ge1-\sum_{j=1}^ke_j\ge\frac12.
\]
The induction closes. Finally,
\[
  \Bigl|\frac{N_{n,m}}{M_n}-1\Bigr|
  =\Bigl|\prod_{j=1}^{n-1}(1+\xi_j)-1\Bigr|
  \le e^{\sum_j|\xi_j|}-1\le e^{E(n,t)}-1\le\varepsilon(n,t).
\]
Since \(M_n=B_{n,4t}\), the proof is complete.
\end{proof}

\section{Asymptotic consequences}\label{sec:asymptotic}

\begin{proof}[Proof of Theorem~\ref{thm:uniform}]
If \(\Theta=m/(n^2L)\to\infty\), then eventually \(\Theta\ge40000\). By
Lemmas~\ref{lem:controlE} and~\ref{lem:controlS},
\(E(n,t)+S(n,t)\le F(\Theta)+G(\Theta)\to0\). Thus
Theorem~\ref{thm:explicit} gives the claimed uniform convergence.
\end{proof}

\begin{lemma}[Power laws imply the natural regime]\label{lem:powerlaw}
Let \(0<\delta\le1/2\), put \(u=2+2\delta\) and \(K=6+2u/\delta\). If
\(\omega:=t/n^u\ge1\) and \(n\ge2\), then
\[
  \frac m{n^2\log(2\pi m)}\ge\frac2K\omega^{1/2}.
\]
\end{lemma}

\begin{proof}
Since \(m=4n^u\omega\),
\(L=\log(8\pi n^u\omega)=\log(8\pi)+u\log n+\log\omega\). Using
\(\log x\le x^a/a\) for \(x\ge1\), first with \(a=\delta/2\) and then with
\(a=1/2\), gives \(\log n\le(2/\delta)n^{\delta/2}\) and
\(\log\omega\le2\omega^{1/2}\). Since \(\log(8\pi)\le6\) and
\(n^{\delta/2},\omega^{1/2}\ge1\), we get
\(L\le K(n^{\delta/2}+\omega^{1/2})\). Therefore, writing
\(n^u=n^2n^{2\delta}\),
\[
  \frac m{n^2L}\ge\frac{4n^{2\delta}\omega}{K(n^{\delta/2}+\omega^{1/2})}.
\]
Finally \(n^{\delta/2}\le n^{2\delta}\) and \(\omega^{1/2}\le\omega\), so
\(\omega^{1/2}(n^{\delta/2}+\omega^{1/2})\le2n^{2\delta}\omega\), which is the
claim.
\end{proof}

\begin{proof}[Proof of Corollary~\ref{cor:exponent}]
Given an exponent \(u>2\), put \(\delta=\min\{(u-2)/2,1/2\}\) and
\(u_0=2+2\delta\le u\). The assumption \(t/n^u\to\infty\) implies
\(t/n^{u_0}\to\infty\). Hence Lemma~\ref{lem:powerlaw} shows that
\(m/(n^2L)\to\infty\). Apply Theorem~\ref{thm:uniform}.
\end{proof}

Lemma~\ref{lem:powerlaw} is sharp, in the following precise sense.

\begin{lemma}[No exponent at most two suffices]\label{lem:negative}
Let \(1\le u\le2\). For all \(N\ge2\) and \(\omega\ge1\), writing
\(m=4N^u\omega\),
\begin{equation}\label{eq:negative}
  \frac m{N^2\log(2\pi m)}\le\frac{4\omega}{\log N}.
\end{equation}
Consequently there is a family with \(\omega\to\infty\) along which the left
side stays bounded by \(1/2\): take \(\omega=j\) and \(N=e^{8j}\), \(j\ge1\).
\end{lemma}

\begin{proof}
Since \(u\le2\) and \(N\ge2\), we have \(N^u\le N^2\); and since \(u\ge1\),
\(\omega\ge1\), \(\log(2\pi m)=\log(8\pi)+u\log N+\log\omega\ge\log N>0\).
Hence
\[
  \frac m{N^2\log(2\pi m)}=\frac{4N^u\omega}{N^2\log(2\pi m)}
  \le\frac{4N^2\omega}{N^2\log N}=\frac{4\omega}{\log N}.
\]
For the family, \(\log N=8j=8\omega\), so the bound is \(4\omega/(8\omega)=1/2\).
\end{proof}

The witness can be taken to consist of genuine \emph{integer} pairs, which is
the domain of the counting problem; the relaxation to real parameters above is
a convenience, not a load-bearing step.

\begin{lemma}[Integer witness]\label{lem:intwitness}
Let \(1\le u\le2\) and let \(j\ge1\) be an integer. Put
\[
  n_j:=\bigl\lceil e^{16j}\bigr\rceil,\qquad
  t_j:=\bigl\lceil j\,n_j^u\bigr\rceil,\qquad m_j:=4t_j .
\]
Then \(n_j\ge2\), \(t_j\ge1\), \(t_j/n_j^u\ge j\), and
\[
  \frac{m_j}{n_j^2\log(2\pi m_j)}\le\frac12 .
\]
\end{lemma}

\begin{proof}
Clearly \(n_j\ge2\) and \(t_j\ge1\), and \(t_j\ge j\,n_j^u\) gives
\(t_j/n_j^u\ge j\), which tends to infinity. Ceilings cost at most one, so
\(t_j\le j\,n_j^u+1\); since \(u\le2\) and \(4\le4j\,n_j^2\),
\[
  m_j=4t_j\le4j\,n_j^u+4\le4j\,n_j^2+4j\,n_j^2=8j\,n_j^2 .
\]
Since \(u\ge1\) and \(j\ge1\) we have \(n_j\le n_j^u\le j\,n_j^u\le t_j\),
whence \(2\pi m_j\ge n_j\) and \(\log(2\pi m_j)\ge\log n_j\ge16j\). Therefore
\[
  \frac{m_j}{n_j^2\log(2\pi m_j)}\le\frac{8j\,n_j^2}{n_j^2\cdot16j}=\frac12 .
  \qedhere
\]
\end{proof}

The family of Lemma~\ref{lem:negative} satisfies the power law with exponent
\(u\) with room to spare, \(t/N^u=\omega=j\to\infty\), yet the
near-quadratic quantity never exceeds \(1/2\). So no exponent \(u\le2\) forces
the hypothesis of Theorem~\ref{thm:uniform}, while by
Lemma~\ref{lem:powerlaw} every \(u>2\) does. The infimum of the exponents
reachable by this argument is therefore exactly \(2\), and it is not attained.
Both directions are formally verified; in the accompanying Lean development the
positive direction is \texttt{bridges\_of\_two\_lt}, the negative direction
\texttt{not\_forces\_of\_le\_two} with its integer form
\texttt{exists\_int\_witness\_of\_le\_two}, and the conclusion
\texttt{sInf\_admissibleExponents}.

A remark on what ``forces'' must mean here. The relevant notion is the one used
in Corollary~\ref{cor:exponent}: along every sequence of pairs \((n_i,t_i)\)
with \(\omega_i=t_i/n_i^u\to\infty\), and with \(n_i\) free to vary, the
quantity \(m_i/(n_i^2\log2\pi m_i)\) tends to infinity. The uniformity in \(n\)
is essential, not a technical convenience: if instead \(n\) were held fixed
while \(\omega\to\infty\), then \(m/(n^2\log2\pi m)\) would diverge like
\(\omega/\log\omega\) for \emph{every} exponent \(u\), so every \(\alpha\ge1\)
would count as admissible and the constant would carry no information. In the
Lean development this equivalence is
\texttt{forces\_iff\_sequential}, which shows the quantifier form used there
agrees with the sequential premise above.

We stress what the negative direction does and does not say. It shows that the
\emph{implication} ``power law with exponent \(u\le2\) \(\Rightarrow\)
near-quadratic hypothesis'' is false, which is precisely why the present method
stops at \(2\). It does \emph{not} show that the asymptotic formula itself
fails for such \(u\); whether it holds there is open, as recorded in
Section~\ref{sec:limits}.

\section{Exact low-dimensional checks}

The correction factor can be checked against exact formulas. For
\(m\in4\N\),
\[
  N_{2,m}=2^m\binom m{m/2},\qquad
  N_{3,m}=2^m\frac{m!}{((m/4)!)^4},\qquad
  N_{4,m}=2^m\sum_{a+b=m/4}\frac{m!}{(a!)^4(b!)^4}.
\]
For the last identity, normalize the first row to all \(+1\). The eight column
types of the remaining three rows have multiplicities determined by a single
third-order Fourier coefficient: the four types with product \(+1\) occur
\(a\) times each and the four types with product \(-1\) occur \(b\) times each.

\begin{table}[ht]
\caption{Relative error with and without the correction factor.}
\begin{tabular}{ccrr}
\toprule
\(n\) & \(m\) & \(N_{n,m}/B_{n,m}-1\) & \(N_{n,m}/A_{n,m}-1\)\\
\midrule
2 & 8   & \(+7.99\times10^{-5}\) & \(-3.069\times10^{-2}\)\\
2 & 32  & \(+1.27\times10^{-6}\) & \(-7.781\times10^{-3}\)\\
2 & 256 & \(+2.48\times10^{-9}\) & \(-9.761\times10^{-4}\)\\
3 & 8   & \(+1.30\times10^{-3}\) & \(-1.435\times10^{-1}\)\\
3 & 32  & \(+2.15\times10^{-5}\) & \(-3.829\times10^{-2}\)\\
3 & 256 & \(+4.22\times10^{-8}\) & \(-4.871\times10^{-3}\)\\
4 & 16  & \(+1.49\times10^{-3}\) & \(-1.953\times10^{-1}\)\\
4 & 32  & \(-3.76\times10^{-4}\) & \(-1.039\times10^{-1}\)\\
4 & 40  & \(-2.61\times10^{-4}\) & \(-8.402\times10^{-2}\)\\
\bottomrule
\end{tabular}
\end{table}

These calculations are consistency checks, not inputs to the proof.

\section{Scope and limitations}\label{sec:limits}

The local expansion in Proposition~\ref{prop:central} has accumulated relative
error of order \(n^4/m^2\), which already tends to zero under
\(m/n^2\to\infty\). The logarithmic loss comes entirely from comparing the
exceptional set of far-field matrices, whose ambient density is
\(e^{-\Omega(m)}\), with the density of partial Hadamard matrices, which is
roughly \(e^{-\Theta(n^2\log m)}\). Removing the logarithm would require new
information about the far field under the partial-Hadamard conditioning; a
universal pointwise estimate is ruled out by Remark~\ref{rem:exceptional}.

The proof establishes a corrected asymptotic formula. It does not imply
\(N_{n,m}\sim A_{n,m}\) below the cubic scale, because \(R_{n,m}\) can tend to
zero. Nor does it prove the endpoint regime \(m/n^2\to\infty\). In particular,
the upper bound \(2\) for the critical exponent is an infimum over admissible
power laws and is not attained by the present argument.

\section{Use of AI tools}

The proof strategy and an initial detailed draft were developed with the
assistance of Hyra, a research agent \cite{hyra2026}. The present manuscript
was subsequently subjected to an independent line-by-line mathematical audit,
and local issues in the original draft were repaired, including the description
of the universal saddle lattice and an endpoint term in the net-cardinality
estimate. AI-generated judgments were not treated as certificates of
correctness; the mathematical claims rest on the explicit argument above. This
disclosure follows the general principle used in other Hyra-assisted
mathematical work \cite{LinLi2026}.

\appendix
\section{Numerical constants}

For convenience, we collect the nontrivial numerical inequalities used in the
proof:
\[
  \log(8\pi)>3.2,\qquad
  -\log(\cos1)-\tfrac7{12}<\tfrac1{30},\qquad
  e^{1/24}<1.043,\qquad
  \sqrt{64/3}<e^{1.54},
\]
\[
  \tfrac3{256}>\tfrac1{86},\qquad
  \tfrac14+\tfrac1{38}<0.28,\qquad
  \log200<1.66\cdot3.2,\qquad
  2\cdot5.66\cdot1024<20000.
\]
All remaining constants are exact rational quantities or are obtained by
rounding in the safe direction.

\bibliographystyle{amsplain}
\bibliography{reference}

\end{document}
